
\phantomsection\label{thm:sga3-V-9.0}%
\noindent\textbf{Theorem 9.0.}---%
\textit{Let \(S\) be an arbitrary scheme and}
\begin{equation*}
  \resizebox{0.34\textwidth}{!}{%
    \SGAThreeNativeInput{figures/exp5/exp5_localp032_theorem90_groupoid.png}}
  \phantomsection\label{fig:sga3-V-9.0-groupoid}
\end{equation*}
\textit{a \((\mathbf{Sch}/S)\)-groupoid such that \(X_0\) is of
finite presentation and quasi-projective over \(S\), \(d_1\) is of
finite presentation, proper, and flat, and the morphism}
\[
  d_0\boxtimes d_1:X_1\longrightarrow X_0\times_S X_0
\]
\textit{is quasi-finite. Then:}
\begin{enumerate}
\item[(1)]\phantomsection\label{thm:sga3-V-9.0-1}%
\textit{Every point \(x\) of \(X_0\) has a saturated open
neighborhood \(W\) such that the groupoid induced by \(X_\ast\) on
\(W\) has a quasi-section.}

\item[(2)]\phantomsection\label{thm:sga3-V-9.0-2}%
\textit{Let \((Y,p)\) be the cokernel of \((d_0,d_1)\) in the
category of all ringed spaces. Then \(Y\) is a scheme, \(p\) is a
morphism of schemes, and \((Y,p)\) is a cokernel of \((d_0,d_1)\) in
\((\mathbf{Sch}/S)\).}

\item[(3)]\phantomsection\label{thm:sga3-V-9.0-3}%
\textit{The morphism \(p\) is surjective, open, and universally
closed.}

\item[(4)]\phantomsection\label{thm:sga3-V-9.0-4}%
\textit{The morphism}
\[
  d:X_1\longrightarrow X_0\times_Y X_0
\]
\textit{with components \(d_0\) and \(d_1\) is surjective.}

\item[(5)]\phantomsection\label{thm:sga3-V-9.0-5}%
\textit{If \((d_0,d_1)\) is an equivalence pair, then:}
  \begin{enumerate}
  \item[(a)]\phantomsection\label{thm:sga3-V-9.0-5-a}%
  \textit{the morphism}
  \[
    d:X_1\longrightarrow X_0\times_Y X_0
  \]
  \textit{is an isomorphism and \(p\) is faithfully flat;}

  \item[(b)]\phantomsection\label{thm:sga3-V-9.0-5-b}%
  \textit{the morphisms \(p\) and \(Y\to S\) are of finite
  presentation, and \((Y,p)\) is a cokernel of \((d_0,d_1)\) in the
  category of sheaves for the fppf topology.}
  \end{enumerate}
\end{enumerate}

\noindent\textit{Proof.} For \hyperref[thm:sga3-V-9.0-1]{(1)}, the question is local on \(S\), so
we may suppose that \(S=\operatorname{Spec}B\) is affine. There then
exist a ring \(A\) of finite type over \(\mathbb Z\), a morphism
\[
  S\longrightarrow T=\operatorname{Spec}A,
\]
and a \((\mathbf{Sch}/T)\)-groupoid \(Z_\ast\) such that
\[
  X_\ast\simeq Z_\ast\times_T S.
\]
Indeed, apply EGA~IV\textsubscript{3}, 8.8.3 to
\(S_0=\operatorname{Spec}\mathbb Z\) and
\(S_i=\operatorname{Spec}A_i\), where the \(A_i\) range over the
finitely generated \(\mathbb Z\)-subalgebras of \(B\). Moreover, by
EGA~IV\textsubscript{3}, 8.10.5, we may suppose that \(Z_\ast\)
satisfies the hypotheses of
Theorem~\hyperref[thm:sga3-V-7.1]{7.1}. Consequently, \(Z_\ast\) has
quasi-sections locally.

By part~\hyperref[par:sga3-V-9-a]{(a)}, the same is therefore true of
\(X_\ast\), and assertions \hyperref[thm:sga3-V-9.0-2]{(2)}, \hyperref[thm:sga3-V-9.0-3]{(3)}, \hyperref[thm:sga3-V-9.0-4]{(4)}, and \hyperref[thm:sga3-V-9.0-5-a]{(5)(a)} follow from
Lemma~\hyperref[lem:sga3-V-6.1]{6.1}, just as in the proof of
Theorem~\hyperref[thm:sga3-V-7.1]{7.1}.

\medskip
\phantomsection\label{proof:sga3-V-9.0-c}%
\noindent\textbf{c) Proof of \hyperref[thm:sga3-V-9.0-5-b]{(5)(b)}.}---%
Let us show that \(Y\to S\) is of finite
presentation.\footnote{\label{note:sga3-V-47}Editor's note: The
original states that this follows from
Proposition~\hyperref[sga3:v:proposition:9:1]{9.1} below. We were unable to
reconstruct that argument. The proof that follows was communicated
to us by O.~Gabber.}
By hypothesis, \((d_0^{X_\ast},d_1^{X_\ast})\) is an equivalence
pair; that is,
\[
  d^{X_\ast}:X_1\longrightarrow X_0\times_S X_0
\]
is a monomorphism. By EGA~IV\textsubscript{3}, 8.10.5, after enlarging
\(A\) if necessary, we may suppose that
\[
  d^{Z_\ast}:Z_1\longrightarrow Z_0\times_T Z_0
\]
is a monomorphism. Since \(T=\operatorname{Spec}A\), with \(A\)
Noetherian, Theorem~\hyperref[thm:sga3-V-7.1]{7.1} then shows that the
groupoid \(Z_\ast\) has a cokernel \((Q,q)\) in
\((\mathbf{Sch}/T)\), that \(q\) and \(Q\to T\) are of finite
presentation, and, moreover, that
\[
  q:Z_0\longrightarrow Q
\]
is faithfully flat and \(d^{Z_\ast}\) induces an isomorphism
\[
  Z_1\xrightarrow{\ \sim\ }Z_0\times_Q Z_0.
\]
Put \(Q_S=Q\times_T S\). Since \(X_i\simeq Z_i\times_T S\), we obtain
an isomorphism
\[
  d^{Z_\ast}\times_T S:
  X_1\xrightarrow{\ \sim\ }X_0\times_{Q_S}X_0.
\]
Denote its inverse by \(\phi\), and let \(\pi\) be the canonical
morphism
\[
  X_0\times_Y X_0\longrightarrow X_0\times_{Q_S}X_0.
\]

Then \(\phi\circ\pi\) is the inverse of
\[
  d_0\boxtimes d_1:
  X_1\xrightarrow{\ \sim\ }X_0\times_Y X_0.
\]
It follows that the equivalence relation defined by \(X_\ast\), that
is, the monomorphism shown in
Figure~\hyperref[fig:sga3-V-9.0-equivalence-relation]{\ref*{fig:sga3-V-9.0-equivalence-relation}},
\SGAThreeDiagram[0.78\textwidth]
  {figures/exp5/exp5_localp033_equivalence_relation.png}
  {The equivalence-relation monomorphism defined by \(X_\ast\).}
  {fig:sga3-V-9.0-equivalence-relation}
  {Expos\'e V, component-local p.~33 / combined-reader p.~295 / re-edited header p.~283}
identifies with the equivalence relation \(\mathscr R(q_S)\) defined
by \(q_S:X_0\to Q_S\). Since \(q_S\) is faithfully flat and of finite
presentation, hence a universal effective epimorphism,
\(\mathscr R(q_S)\) has quotient \(Q_S\); see \SGARefLink{sga3:iv:definition:3:3:2}{Expos\'e IV, 3.3.2}.
Consequently,
\[
  Y\simeq Q\times_T S,
\]
so \(Y\to S\) and \(p\) are of finite presentation. Moreover, by
\SGARefLink{sga3:iv:corollary:6:3:3}{Expos\'e IV, 6.3.3}, \((Y,p)\) is also a cokernel of \((d_0,d_1)\) in
the category of sheaves for the fppf topology.

\medskip
\SGARefTarget{sga3:v:proposition:9:1}
\phantomsection\label{prop:sga3-V-9.1}%
\noindent\textbf{Proposition 9.1.}---%
\textit{Consider morphisms of schemes}
\[
  X_0\xrightarrow{\,p\,}Y\xrightarrow{\,q\,}S
\]
\textit{such that \(qp\) is of finite type (respectively, of finite
presentation) and \(p\) is faithfully flat and of finite
presentation. Then \(q\) is of finite type (respectively, of finite
presentation).}%
\begingroup
\renewcommand{\thefootnote}{*}%
\footnote{\label{note:sga3-V-9.1-generalization}See
EGA~IV\textsubscript{4}, 17.7.5 for a more general result.}
\endgroup
\addtocounter{footnote}{-1}%

\noindent Since \(p\) is surjective and \(qp\) is quasi-compact,
\(q\) is quasi-compact. We may therefore suppose that \(S\), \(Y\),
and \(X_0\) are affine, with rings \(A\), \(B\), and \(C\),
respectively. We have
\[
  B=\varinjlim_i B_i,
\]
where the \(B_i\) range over the finitely generated
\(A\)-subalgebras of \(B\). Since \(C\) is of finite presentation
over \(B\), there exist an index \(i_0\), a \(B_{i_0}\)-algebra
\(C_{i_0}\) of finite presentation, and an isomorphism
\[
  C\simeq C_{i_0}\otimes_{B_{i_0}}B.
\]
For \(i\geq i_0\), put
\[
  C_i=C_{i_0}\otimes_{B_{i_0}}B_i;
\]
then
\[
  C\simeq C_i\otimes_{B_i}B.
\]

\SGAThreeDiagram[0.22\textwidth]
  {figures/exp5/exp5_localp033_algebra_limit_diagram.png}
  {The compatible \(A\)-algebras used in the limit argument.}
  {fig:sga3-V-9.1-algebra-limit}
  {Expos\'e V, component-local p.~33 / combined-reader p.~295 / re-edited header p.~283}

Since \(C\) is faithfully flat over \(B\), EGA~IV\textsubscript{3},
11.2.6 and 8.10.5(vi) give an index \(i_1\geq i_0\) such that
\(C_{i_1}\) is faithfully flat over \(B_{i_1}\). Consequently,
\(C_i\) is faithfully flat over \(B_i\) for every \(i\geq i_1\).
For \(i\geq i_1\), the canonical map \(C_i\to C\) is then injective,
because it is obtained from \(B_i\to B\) by faithfully flat base
extension.


If \(C\) is of finite type over \(A\), it follows that there is an
index \(j\geq i_1\) such that \(C_j=C\), whence \(B_j=B\), because
\(C_j\) is faithfully flat over \(B_j\). Thus \(B\) is of finite type
over \(A\).

Now suppose that \(C\) is of finite presentation over \(A\). By what
precedes, \(B\) is of finite type over \(A\), and hence has the form
\[
  B=\overline B/I,
\]
where \(\overline B\) is a polynomial algebra over \(A\) in finitely
many indeterminates and \(I\) is an ideal of \(\overline B\). The
ideal \(I\) is the union of its finitely generated subideals
\(I_\alpha\); consequently,
\[
  B=\varinjlim_\alpha B_\alpha,
  \qquad B_\alpha=\overline B/I_\alpha.
\]
Proceeding as above, there exist an index \(\alpha_0\), a
\(B_{\alpha_0}\)-algebra \(C_{\alpha_0}\) of finite presentation,
and an isomorphism
\[
  C\simeq C_{\alpha_0}\otimes_{B_{\alpha_0}}B.
\]
For \(\alpha\geq\alpha_0\), put
\[
  C_\alpha=C_{\alpha_0}\otimes_{B_{\alpha_0}}B_\alpha,
\]
so that
\[
  C\simeq C_\alpha\otimes_{B_\alpha}B
  \qquad(\alpha\geq\alpha_0).
\]
Again by EGA~IV\textsubscript{3}, 11.2.6 and 8.10.5(vi),
\(C_\alpha\) is faithfully flat over \(B_\alpha\) for all sufficiently
large \(\alpha\). In that case, the kernels of the maps
\(C_\alpha\to C\) and \(C_\alpha\to C_\beta\), respectively, identify
with
\[
  C_\alpha\otimes_{B_\alpha}(I/I_\alpha)
  \quad\text{and}\quad
  C_\alpha\otimes_{B_\alpha}(I_\beta/I_\alpha)
  \qquad(\beta\geq\alpha).
\]

Since \(C_\alpha\) and \(C\) are of finite presentation over \(A\)
and \(C_\alpha\to C\) is surjective,
\[
  C_\alpha\otimes_{B_\alpha}(I/I_\alpha)
\]
is a finitely generated ideal,\footnote{\label{note:sga3-V-48}Editor's
note: see EGA~IV\textsubscript{1}, 1.4.4.} and it is the union of the
ideals
\[
  C_\alpha\otimes_{B_\alpha}(I_\beta/I_\alpha).
\]
Thus, for all sufficiently large \(\beta\),
\[
  C_\alpha\otimes_{B_\alpha}(I_\beta/I_\alpha)
  =
  C_\alpha\otimes_{B_\alpha}(I/I_\alpha).
\]
Since \(C_\alpha\) is faithfully flat over \(B_\alpha\), this also
gives \(I_\beta=I\). Therefore \(B\) is of finite presentation over
\(A\).

\bigskip
\phantomsection\label{sec:sga3-V-10}%
\section*{10. Supplement: quotients by a group scheme}
\addcontentsline{toc}{section}{10. Supplement: quotients by a group scheme}

Sections \hyperref[sec:sga3-V-10.2]{10.2}--\hyperref[sec:sga3-V-10.4]{10.4} below, written following suggestions of
M.~Raynaud, apply the preceding theorems to the case of an action of
a group scheme. For the reader's convenience, \hyperref[sga3:v:paragraph:10-1]{Section~10.1} first
reproduces statements 2.1--2.3 of Expos\'e XVI.

\medskip
\phantomsection\label{subsec:sga3-V-10.1}%
\noindent\SGAInlineRefTarget{sga3:v:paragraph:10-1}\textbf{10.1. Representability theorems for quotients.}---%
We first recall \hyperref[thm:sga3-V-10.1.1]{the following result}.

\medskip
\phantomsection\label{thm:sga3-V-10.1.1}%
\noindent\textbf{Theorem 10.1.1.}---%
\textit{Let \(S\) be a scheme, let \(X\) and \(Y\) be two
\(S\)-schemes, and let \(f:X\to Y\) be an \(S\)-morphism. Suppose
that one of the following two conditions holds:}
\begin{enumerate}
\item[\(\alpha\))]\phantomsection\label{thm:sga3-V-10.1.1-alpha}%
\textit{the morphism \(f\) is locally of finite presentation;}
\item[\(\beta\))]\phantomsection\label{thm:sga3-V-10.1.1-beta}%
\textit{the scheme \(S\) is locally Noetherian and \(X\) is locally
of finite type over \(S\).}
\end{enumerate}
\textit{Then the following conditions are equivalent:}
\begin{enumerate}
\item[(i)]\phantomsection\label{thm:sga3-V-10.1.1-i}%
\textit{there exist an \(S\)-scheme \(X'\) and a factorization}
\[
  f:X\xrightarrow{\,f'\,}X'\xrightarrow{\,f''\,}Y,
\]
\textit{where \(f'\) is a faithfully flat \(S\)-morphism locally of
finite presentation and \(f''\) is a monomorphism;}

\item[(ii)]\phantomsection\label{thm:sga3-V-10.1.1-ii}%
\textit{the first projection}
\[
  p_1:X\times_Y X\longrightarrow X
\]
\textit{is flat.}
\end{enumerate}
\textit{Moreover, if these conditions hold, then \((X',f')\) is a
quotient of \(X\) by the equivalence relation defined by \(f\) for
the fppf topology. Thus the factorization \(f=f''\circ f'\) in \hyperref[thm:sga3-V-10.1.1-i]{(i)} is
unique up to isomorphism.}

The case in which \(Y\) is locally Noetherian and \(X\) is of finite
type over \(Y\) is treated in \hyperref[bib:sga3-V-Mur65]{[Mur65]}, Corollary~2 to Theorem~2. We
shall see that the general case reduces to that one.

We first make a few remarks.

\medskip
\phantomsection\label{par:sga3-V-10.1.1-a}%
\noindent\textbf{a)} The implication \hyperref[thm:sga3-V-10.1.1-i]{(i)}\(\Rightarrow\)\hyperref[thm:sga3-V-10.1.1-ii]{(ii)} is
immediate. Indeed, the first projection
\[
  p'_1:X\times_{X'}X\longrightarrow X
\]
factors through \(X\times_Y X\):
\[
  p'_1:
  X\times_{X'}X\xrightarrow{\,u\,}X\times_Y X
  \xrightarrow{\,p_1\,}X.
\]
The morphism \(u\) is an isomorphism because \(f''\) is a
monomorphism, and \(p'_1\) is flat because \(f'\) is flat; hence
\(p_1\) is flat.

\medskip
\phantomsection\label{par:sga3-V-10.1.1-b}%
\noindent\textbf{b)} The assertions of
Theorem~\hyperref[thm:sga3-V-10.1.1]{10.1.1} are local on \(Y\), and
hence local on \(S\). They are also local on \(X\), as follows readily
from the fact that a flat morphism locally of finite presentation is
open (EGA~IV\textsubscript{3}, 11.3.1).

\medskip
\phantomsection\label{par:sga3-V-10.1.1-c}%
\noindent\textbf{c)} Under hypothesis
\(\hyperref[thm:sga3-V-10.1.1-alpha]{\alpha}\) of
Theorem~\hyperref[thm:sga3-V-10.1.1]{10.1.1}, the preceding remarks
reduce us to the case where \(X\) and \(Y\) are affine and \(f\) is
of finite presentation. Replacing \(S\) by \(Y\), if necessary, we
may suppose that \(X\) and \(Y\) are of finite presentation over
\(S\). EGA~IV\textsubscript{3}, 11.2.6 then reduces us to the case
where \(S\) is Noetherian.

\medskip
\phantomsection\label{par:sga3-V-10.1.1-d}%
\noindent\textbf{d)} Under hypothesis
\(\hyperref[thm:sga3-V-10.1.1-beta]{\beta}\) of
Theorem~\hyperref[thm:sga3-V-10.1.1]{10.1.1}, we may suppose that
\(S\), \(X\), and \(Y\) are affine, that \(S\) is Noetherian, and
that \(X\) is of finite type over \(S\). Write \(Y\) as a filtered
inverse limit of affine schemes \(Y_i\) of finite type over \(S\).
The schemes
\[
  X\times_{Y_i}X
\]
form a filtered decreasing family of closed subschemes of
\(X\times_S X\), whose inverse limit is \(X\times_Y X\). Since
\(X\times_S X\) is Noetherian, for all sufficiently large \(i\) we
have
\[
  X\times_{Y_i}X=X\times_Y X.
\]
It follows that
\[
  f_i:X\xrightarrow{\,f\,}Y\longrightarrow Y_i
\]
satisfies \hyperref[thm:sga3-V-10.1.1-ii]{condition (ii)} of
Theorem~\hyperref[thm:sga3-V-10.1.1]{10.1.1} whenever \(f\) does.
The equivalence relation on \(X\) defined by \(f\) is the same as the
one defined by \(f_i\). Thus it suffices to prove
\hyperref[thm:sga3-V-10.1.1-ii]{(ii)}\(\Rightarrow\)\hyperref[thm:sga3-V-10.1.1-i]{(i)} for \(f_i\), reducing us to the case where
\(Y\) is of finite type over \(S\).

\medskip
\noindent\textit{Application to group schemes.}
Let \(S\) be a scheme, let \(G\) be an \(S\)-group scheme locally of
finite presentation over \(S\), and suppose that \(G\) acts on the
left on an \(S\)-scheme \(X\). If \(X\to S\) has a section \(\xi\),
recall that the stabilizer \(\operatorname{Stab}_G(\xi)\) is
representable by a subgroup scheme of \(G\); see \SGARefLink{sga3:i:definition:2:3:3}{Expos\'e I, 2.3.3}.

\medskip
\phantomsection\label{thm:sga3-V-10.1.2}%
\noindent\textbf{Theorem 10.1.2.}---%
\textit{Let \(S\) be a scheme and let \(G\) be an \(S\)-group scheme
locally of finite presentation over \(S\), acting on an \(S\)-scheme
\(X\). Suppose that \(X\to S\) has a section \(\xi\) whose stabilizer
\(H\) in \(G\) is flat over \(S\). If one of the following
hypotheses holds:}
\begin{enumerate}
\item[\textit{a)}]\phantomsection\label{thm:sga3-V-10.1.2-a}%
\textit{\(X\) is locally of finite type over \(S\);}
\item[\textit{b)}]\phantomsection\label{thm:sga3-V-10.1.2-b}%
\textit{\(S\) is locally Noetherian,}
\end{enumerate}
\textit{then the fppf quotient sheaf \(G/H\) is representable by an
\(S\)-scheme, locally of finite presentation over \(S\), and the
\(S\)-morphism}
\[
  f:G\longrightarrow X,\qquad g\longmapsto g\cdot\xi
\]

